\chapter{The Corridor Rotates}

\section{Rotation turns number into phase}

The electron gave the corridor a stable name. This chapter sets the corridor turning, and finds that rotation converts a
static reading into a phase --- that a number at rest, once rotated, becomes an experiment about angle. The pivot of the
conversion is a single operator, the quarter-turn, whose square is minus one: the imaginary unit read as a rotation. When the
corridor turns through this axis, the reading it carried as a magnitude re-presents as a phase, and the physics that was
statics becomes the physics of interference. The turn is the operation; the phase is what the turn makes of the number.

The mathematics is the identity of $i$ with a rotation by a right angle. Multiplication by $i$ takes the real axis to the
imaginary axis and back to the negative real axis in two steps, $i^2 = -1$; it is the generator of turning in the plane,
$e^{i\theta} = \cos\theta + i\sin\theta$. A static real number $x$ lives on the real axis; rotate the axis through $i$ and $x$
acquires a phase, $x \mapsto x\,e^{i\theta}$, and what was a bare magnitude is now a magnitude with an angle. In the quantum
description this is not a formal trick but the whole content of dynamics: the state advances by a phase, $\psi(t) =
e^{-iHt/\hbar}\psi(0)$, and the generator of that advance is the Hamiltonian standing behind the $-i$. The static reading and
the phase are the same number read before and after the corridor has turned.

The $-i$ axis is where statics becomes kinematics without anything being added. Wick's rotation makes this exact: replace the
time by an imaginary time, $t \to -i\tau$, and the oscillating phase $e^{-iHt/\hbar}$ becomes a decaying weight
$e^{-H\tau/\hbar}$ --- the dynamics of phases turns into the statics of a thermal distribution, and back, by a quarter turn
of the corridor. The two regimes the earlier chapters kept apart --- the static reading and the phase experiment --- are one
reading viewed along two axes ninety degrees apart. Rotating the corridor through $-i$ is the operation that carries the one
into the other. Nothing new enters; the same residue is read once as a magnitude and once as a phase, and the axis of the
reading is all that changed.

This is why the instrument treats phase as a genuine experiment and not a bookkeeping convenience. A phase is measurable ---
it is read as an interference fringe, a shift in where constructive and destructive addition fall --- and it is measurable
precisely because it is the rotated form of a magnitude the instrument already held. The corridor did not acquire a new
quantity when it turned; it re-presented an old one along a new axis, where a bench with two arms can read it. The static
number became a phase experiment because rotation is a real operation on the corridor, and the phase it produces is a real
reading, floored and bracketed like any other.

In the two verbs, rotation is a tange that changes which axis the funge is read along. The residue is the funge, bagged
invariant; the choice of axis --- real or imaginary, static or phase --- is the tange, the contravariant selection of how to
read the bag. Turning the corridor through $-i$ re-tanges the same funge along a rotated axis, and the reading changes from
magnitude to phase without the bag changing at all. The number and the phase are one funge under two tanges ninety degrees
apart. The next section catalogues the faces this turning sweeps through --- charge, mass, current, spin, orientation --- and
insists again that which one appears is fixed by where the corridor is viewed.

\section{Charge, mass, current, spin, orientation}

The turn converts a number into a phase; this section catalogues the faces the turning corridor presents as it sweeps
through its axes, and draws once more the lesson that governs the whole part: the face is fixed by the viewing, not by the
residue. Charge, mass, current, spin, orientation --- these are not five possessions of the corridor but five readings it
gives at five angles, and a change of frame carries one into another exactly as a rotation carries the real axis into the
imaginary. The section is a tour of the faces and a reminder that the tour is a tour of viewpoints.

Take the clearest case, the way a boost mixes the electromagnetic faces. To one observer a static charge presents a pure
electric field; to an observer boosted past it, the same charge presents an electric \emph{and} a magnetic field, because
the field strength $F_{\mu\nu}$ is an antisymmetric tensor whose electric and magnetic parts are frame-dependent
projections, $E_i = F_{0i}$, $B_i = -\tfrac12\epsilon_{ijk}F_{jk}$~\cite{maxwell1865,lorentz1904}. The charge did not sprout
a current; the observer turned, and the turning re-projected the one field tensor onto a differently-tilted split of electric
and magnetic. Current is charge seen from a boosted frame; the magnetic face is the electric face rotated. The five names are
five projections of a small number of invariant tensors, and the boost is the rotation that trades them.

Mass and spin join the catalogue as the two invariants that anchor it, and their role is worth marking because it is the
opposite of the others'. Where charge-versus-current and electric-versus-magnetic are frame-dependent faces that the turning
sweeps between, mass and spin are the Casimirs --- the invariants that label the representation and do \emph{not} change
under the turn~\cite{dirac1928}. The mass $m^2 = p_\mu p^\mu$ and the spin are what every frame agrees on; they are the fixed
axis the other faces rotate around. So the catalogue has two kinds of entry: the invariant anchors, mass and spin, that the
turning leaves alone, and the frame-dependent faces, charge-current and electric-magnetic and orientation, that the turning
sweeps through. Both are real; only the second kind depends on the viewing.

This split --- invariant anchor against frame-dependent face --- is the same distinction the whole book has kept between the
residue and its readings, now seen inside a single chapter. The residue's invariants (mass, spin, the conserved charge) are
the funge every frame shares; the residue's faces (the electric-magnetic split, the charge-current split, the orientation) are
the tanges each frame selects. A reader who takes the frame-dependent faces for possessions of the corridor makes the
error the book has warned against since Chapter 3: mistaking a viewpoint for a thing. The corridor possesses its invariants;
it \emph{presents} its faces; and the turning of this chapter is the operation that carries one presentation into another
while leaving the possessions fixed.

In the two verbs, the catalogue is a list of tanges over a fixed funge. The corridor's invariants --- mass, spin, charge ---
are the covariant funge, the bag every frame shares. Its faces --- current, the magnetic side, orientation --- are the
contravariant tanges, the selections a frame makes from the bag as it turns. To sweep the faces is to sweep the tanges while
the funge holds still. So the five names of this section are one invariant residue read at five angles, and the reader who
keeps the funge and the tange apart will never confuse the corridor's possessions with its presentations. The last section
reads the turning itself as a physical quantity --- angular momentum --- and finds that carrying one face into another has a
price the instrument is charged.

\section{Relabeling is angular momentum}

The corridor turns, and turning sweeps its faces. This closing section reads the turning itself, and finds it is not free:
to carry one face into another is to rotate the corridor, and rotation is generated by angular momentum, so every relabeling
is charged a quantum of it. The turn has a price, the price is angular momentum, and the price is the deepest reason the
faces are related quantities rather than independent ones --- they are connected by a rotation, and the rotation costs. This
is where the residue's turning becomes a coupling, and where the book first stands at the site the coupling will later be
read --- and, as always, withholds the magnitude.

Angular momentum is the generator of rotation --- that is its definition in this register, not a fact about spinning bodies
but a fact about the algebra of turns. The operators that turn the corridor obey $[J_i, J_j] = i\hbar\,\epsilon_{ijk}J_k$,
and $J$ is precisely what you apply to rotate a state, $\psi \mapsto e^{-i\theta \cdot J/\hbar}\psi$~\cite{pauli1927}. So to
relabel one face into another --- to carry the electric reading into the magnetic, the static into the phase --- is to act
with $J$, and acting with $J$ spends angular momentum. The relabeling is not a free change of notation; it is a physical
rotation of the corridor, and the corridor charges for it in the currency of turns. Relabeling is angular momentum because
relabeling is rotation and rotation is what angular momentum generates.

The price of a turn is exhibited most sharply in the precession a state suffers when carried around a closed circuit of
turns. Carry a spin around a loop in the space of frames and it returns not where it started but precessed, by an angle that
is the residue of the loop --- Thomas's relativistic precession, the rotation a state accumulates when its sequence of
boosts fails to commute~\cite{goldstein2002}. The precession is the holonomy of the turning, the angular-momentum residue of
a loop in frame space, and it is exactly the kind of returned-altered reading Part II built. Two boosts in different orders
leave the corridor rotated relative to itself; that rotation is a real, measurable precession; and it is the price the loop
of turns was charged. The turn has a residue, and the residue is an angle the bench can read.

Here the book stands, for the first time, at the site where the coupling will be read --- and does with it exactly what it
has done at every such edge. The precession of a bound state's spin against its orbit is the physical home of the coupling:
it is where the residue of the turning shows up as a fine splitting of the state's readings, a splitting whose scale is set
by the very coupling the instrument was built to reach. The book names this site --- the spin-orbit precession, the fine
structure of the readings --- and it withholds the scale. It says \emph{where} the coupling lives: in the price of the turn,
the angular-momentum residue of the loop of boosts. It does not say \emph{how much} the price is. The magnitude is the last
part's, and even there it will arrive as a bracket, floored, never as a value read off here. The site is named; the number is
kept.

In the two verbs, relabeling-as-angular-momentum is the tange that turns out to be a funge in disguise --- a selection that
cannot be made without paying the world. To relabel a face is a tange, a change of how we read the bag; but the turning that
performs it is a covariant rotation of the corridor, a funge the world charges in angular momentum. So the contravariant
relabeling and the covariant price are the two sides of one operation: our change of reading is the world's rotation, and the
world bills the turn. The corridor rotates, and every rotation is angular momentum spent. Part III ends with the faces fully
mapped and the turning fully priced --- and the instrument standing at the coupling's site, the price of the turn named and
the number withheld. The next part brings the residue to the physical bench, where the price is put under an apparatus and
made to show.
